\documentclass[10pt,a4paper]{amsart}
\usepackage[margin=19mm]{geometry}
\usepackage{amsmath,amssymb,mathtools}
\usepackage[scr=rsfs,cal=euler]{mathalfa}
\usepackage{xfrac}
\usepackage[unicode,hidelinks,bookmarks=false]{hyperref}
\newcommand{\ie}{i.e.\ }
\newcommand{\cf}{cf.\ }
\newcommand{\resp}{respectively\ }
\newcommand{\Subsection}{\subsection}
\providecommand{\nobreakdash}{\nobreak\hbox{-}}
\setlength{\emergencystretch}{2em}
\multlinegap0pt
\usepackage{bm}
\usepackage{bbm}
\usepackage{dsfont}
\usepackage{xspace}
\usepackage{cancel}
\usepackage{mathtools}
\usepackage{cleveref}
\usepackage{amsthm}
\makeatletter
\@ifundefined{definition}{\newtheorem{definition}{Definition}}{}
\@ifundefined{proposition}{\newtheorem{proposition}{Proposition}}{}
\makeatother
\title[Nonlinear excitations in multi-dimensional nonlocal lattices]{Nonlinear excitations in multi-dimensional nonlocal lattices}
\author{Brian Choi}
\date{Source: arXiv:2408.11177v2 (CC BY 4.0)}
\begin{document}
\maketitle

\noindent\emph{Source and licence.} Nonlinear excitations in multi-dimensional nonlocal lattices. Version arXiv:2408.11177v2 (CC BY 4.0), distributed under the Creative Commons Attribution 4.0 International licence, as recorded in the arXiv licence metadata for this exact version and shown on the abstract page https://arxiv.org/abs/2408.11177v2. Full rights notice: licenses/arxiv-2408.11177-CC-BY-4.0.txt.\par\smallskip

\noindent\textit{Editorial scope.} Counted unit: the Proposition whose source label is \texttt{algebraic\_decay} in the article (source0.tex lines 315--324, source label \texttt{algebraic\_decay}), delivered together with the article's own complete original proof of it (source0.tex lines 325--379, one \texttt{proof} environment of 55 lines). No mathematics is written, repaired or re-derived: statement and proof are the article's own text line for line. Required context retained uncounted: def:kernel (lines 90--104); main\_model3 (lines 197--199); ground\_state (lines 302--304). Every definition, display and label of those lines is retained unchanged. Omitted from this package and not redistributed: 1--89, 105--196, 200--301, 305--314, 380--1020. Nothing inside a retained excerpt is omitted or altered: every retained body is delivered exactly as the article's lines are written, so no omission receipt of any kind accompanies this package. Citations: the retained excerpts carry no citation command at all, so no bibliography entry of the article is transcribed and no reference is redistributed. Reference adaptation: none was needed and none was made; every reference inside the delivered unit resolves inside it, so no unresolved reference or citation is produced. Printed slips kept literal: no printed word, symbol, hypothesis or number of the counted statement or of its proof was altered. Layout and class: the article's own class is \texttt{article} with packages including graphicx, inputenc, authblk, subcaption, comment, amsmath, amsbsy, amssymb, dsfont, upgreek, textcomp, braket, setspace, blindtext; the wrapper is the lane's pinned \texttt{amsart} editorial wrapper at 10pt on A4, which restates the theorem-like environments the retained text needs and adds the article's own macro definitions that the retained text needs (none), with extra wrapper packages (bm, bbm, dsfont, xspace, cancel, mathtools, cleveref). The delivered numbering is the unit's own. Assets: no retained span contains a figure environment, a graphics-inclusion command, a PDF-page inclusion, a table or a TikZ picture; no image file is copied into this package, the package declares no asset dependency and no image file is redistributed with it. Licence: version arXiv:2408.11177v2 of this article is distributed under the Creative Commons Attribution 4.0 International licence, as recorded in the arXiv licence metadata for this exact version and shown on the abstract page https://arxiv.org/abs/2408.11177v2. A full rights notice is delivered at licenses/arxiv-2408.11177-CC-BY-4.0.txt. Characters: every retained excerpt is pure ASCII, so the delivered engine, XeLaTeX with its default Latin Modern text font, reports no missing character.
\begin{definition}\label{def:kernel}

Let $J:\mathbb{Z}^d \rightarrow [0,\infty)$ satisfy $J(0,\dots,0) = 0$ (no self-interaction), $J_{n} = J_{-n}$ for all $n \in \mathbb{Z}^d$ (symmetry), and $\overline{J} := \sum\limits_{m \in \mathbb{Z}^d} J_m < \infty$ (summability). Let $\eta:\mathbb{N} \rightarrow [0,\infty)$ be an envelope such that (i) if $\alpha \in (0,\infty)$, assume there exist $0< C_1 < C_2$ and $k_0 \in \mathbb{N}$ such that $C_1 \leq k^{d+\alpha} \eta(k) \leq C_2$ for all $k \geq k_{0}$; (ii) if $\alpha = \infty$, assume $\lim\limits_{k \rightarrow \infty} k^{d+\beta} \eta(k) = 0$ for any $\beta < \infty$. Further assume 
\begin{enumerate}
    \item[(I)] \textbf{Stable-like in the far-field:} there exist $R, c_{\pm} > 0$ such that for any $k \geq R$ and $n \in \mathbb{Z}^d$ with $|n| \geq R$, 
    \begin{equation*}
    J_n \leq c_{+} \eta(|n|),\qquad \overline{J}_k := \sum_{m \in S_k} J_m \geq c_{-} k^{d-1} \eta(k).    
    \end{equation*} 
    
    \item[(II)] \textbf{Nearest-neighbor nondegeneracy:}
    \begin{equation*}
        \min_{m \in S_1} J_m > 0.
    \end{equation*}    
\end{enumerate}
\end{definition}

    \begin{equation}\label{main_model3}
        -w q_* = L q_* - q_*^p.
    \end{equation}

\begin{equation}\label{ground_state}
    q = (\omega+L)^{-1}q^p = K \ast q^p
\end{equation}

\begin{proposition}\label{algebraic_decay}
Let $q \in l^2(\mathbb{Z}^d)$ satisfy \eqref{main_model3} with $M[q] = \nu$. For any $\alpha \in (0,\infty)$, the decay estimate 
\begin{equation}\label{decay2}
    |q_n| \leq C(d,\alpha,p,\omega) \nu^{\frac{p}{2}} \langle n \rangle^{-(d+\alpha)},\ \forall n \in \mathbb{Z}^d 
\end{equation}
holds, and for $\alpha = \infty$, $q_n = O_{\beta}(|n|^{-\beta})$ for any $\beta > 0$. Furthermore, if $\epsilon > 0,\ \alpha \in (0,\infty)$, then
\begin{equation}
q_n \notin O(|n|^{-(d+\alpha+\epsilon)}).    
\end{equation}
\end{proposition}

\begin{proof}
By an explicit computation of $\mathcal{F}^{-1} [\sin^2 \left(\frac{m \cdot \xi}{2}\right)]_n$, we have
\begin{equation}\label{fourier}
    s_n := [\mathcal{F}^{-1} \sigma]_n = \sum_{m \neq 0} J_{m} \mathcal{F}^{-1} [\sin^2 \left(\frac{m \cdot \xi}{2}\right)]_n =
    \begin{cases}
            \frac{1}{2} \sum\limits_{m \neq 0} J_{m}
            & ,\text{ if } n = 0\\
            -\frac{1}{2} J_{n} & ,\text{ if } n \neq 0.
    \end{cases}
\end{equation}
Assume $\alpha \in (0,\infty)$. To obtain a decay estimate of $q$ from \eqref{ground_state}, an argument similar to the Picard iteration is implemented. Note that the decay of $K$ corresponds to the regularity of $(\omega + \sigma(\xi))^{-1}$, and therefore of $\sigma(\xi)$, which in turn corresponds to the decay of $s_n$; the symbols $(\omega + \sigma(\xi))^{-1}$ and $\sigma(\xi)$ are in the same H\"older space since $\sigma \geq 0$ and $\omega > 0$. In short, $K = O(|n|^{-(d+\alpha)})$.

By $l^2(\mathbb{Z}^d) \hookrightarrow l^\infty(\mathbb{Z}^d)$, for any $\epsilon > 0$, there exists $R>0$ such that $|q_n| \leq \epsilon$ whenever $|n| > R$. For $|n| > 2R$, the convolution in \eqref{ground_state} is split into $|m| \leq R$ and $|m| > R$ to obtain
\begin{equation}\label{ground_state2}
\begin{split}
    q_n &\leq \| q \|_{l^2}^p \sum_{|m| \leq R} K_{n-m} + \epsilon^{p-1}\sum_{|m| > R} K_{n-m} q_m\\
    &\lesssim \| q \|_{l^2}^p\sum_{|m| \leq R} \langle n-m \rangle^{-(d+\alpha)}+ \epsilon^{p-1}\sum_{|m| > R} K_{n-m} q_m =: S_1 + S_2.
\end{split}
\end{equation}
Since $|n-m| \geq \frac{|n|}{2}$, we have
\begin{equation*}
S_1 \leq C  \langle n \rangle^{-(d+\alpha)} \simeq \| q \|_{l^2}^p R^d \langle n \rangle^{-(d+\alpha)}.   
\end{equation*}
For $\delta > 0$, let $C_0(\delta) > 0$ be the best constant that satisfies
\begin{equation}\label{continuity}
    q_n \leq C_0(\delta + \langle n \rangle^{-(d+\alpha)}),\ \forall n \in \mathbb{Z}^d.
\end{equation}
Indeed $C_0$ exists since $\frac{q_n}{\delta + \langle n \rangle^{-(d+\alpha)}} \leq \frac{\nu^{1/2}}{\delta}$. Define $C_1 > 0$ by
\begin{equation*}
\sum\limits_{m} K_{n-m} \langle m \rangle^{-(d+\alpha)} \leq C_1 \langle n \rangle^{-(d+\alpha)},    
\end{equation*}
which can be shown by comparing the series with the corresponding integral. By \eqref{ground_state2}, \eqref{continuity},
\begin{equation*}
\begin{split}
    q_n &\leq C \langle n \rangle^{-(d+\alpha)} + C_0 \epsilon^{p-1} \sum_{|m| > R} K_{n-m} (\delta + \langle m \rangle^{-(d+\alpha)})\\
    &\leq (\epsilon^{p-1}\sum_{m} |K_{m}|) C_0\delta +  (C + C_0 (\epsilon^{p-1}C_1))\langle n \rangle^{-(d+\alpha)}\\
    &\leq (C + \frac{C_0}{2})(\delta + \langle n \rangle^{-(d+\alpha)}),
\end{split}
\end{equation*}
where $0 < \epsilon \ll 1$ was chosen at the last inequality, implying $C_0 \leq C + \frac{C_0}{2}$, and therefore $C_0 \leq 2C$. By \eqref{continuity} and $C>0$, the estimate \eqref{decay2} follows by taking $\delta \rightarrow 0$ in \eqref{continuity}. If $\alpha = \infty$, a similar analysis as above shows $O_{\beta}(|n|^{-\beta})$ decay for any $\beta > 0$.

Suppose there exists $\epsilon>0$ satisfying $d+\alpha + \epsilon \notin \mathbb{Z}$ and $\epsilon < (p-1)(d+\alpha)$ such that $q_n = O(|n|^{-(d+\alpha+\epsilon)})$. By \eqref{ground_state}, this implies $\mathcal{F}^{-1}[\sigma \widehat{q}]_{n} = (Lq)_{n} = (s \ast q)_{n} = O(|n|^{-(d+\alpha+\epsilon)})$. Let $R \gg 1$ such that $R^{-\frac{\epsilon}{2}}\sum\limits_{m} |s_m| \ll 1$ and $|q_{m_0}| > \frac{\| q \|_{l^\infty}}{2}$ for some $|m_0| < R$. Let $2R \leq |n| \leq 2^{-\frac{1}{d+\alpha}} R^{\frac{d+\alpha+\epsilon/2}{d+\alpha}}$ such that
\begin{equation*}
J_{n-m_0} \gtrsim \eta(|n-m_0|) \gtrsim |n|^{-(d+\alpha)},    
\end{equation*}
where the first inequality holds by the averaged lower bound in \Cref{def:kernel} and $|n-m_0| \geq R \gg 1$, and the second inequality holds by the limit property of $\eta$ and $|n-m_0| \leq \frac{3 |n|}{2}$. Then,
\begin{equation}
\begin{split}
    |n|^{-(d+\alpha+\epsilon)}\gtrsim|(s \ast q)_n| &\geq \left|\sum_{|m| \leq R} s_{n-m} q_m\right| - \left|\sum_{|m| > R} s_{n-m} q_m\right|\\
    &\gtrsim \sum_{|m| \leq R} J_{n-m} q_m - \sum_{|m| > R} \frac{|s_{n-m}|}{|m|^{d+\alpha+\epsilon}}\\
    &\gtrsim \| q \|_{l^\infty} |n|^{-(d+\alpha)} - (R^{-\frac{\epsilon}{2}} \sum_{m} |s_m|) R^{-(d+\alpha+\epsilon/2)} \gtrsim |n|^{-(d+\alpha)},
\end{split}
\end{equation}
where an explicit computation \eqref{fourier} and $q_n \geq 0$ were used in the third inequality.    
\end{proof}

\end{document}
